\section*{Introduction générale}

\subsection*{Contexte}

Audit Énergie est une entreprise initialement positionnée sur le courtage en énergie, dont
le métier consiste à accompagner des professionnels dans la renégociation de leurs contrats.
Cette activité lui a donné une connaissance fine des métiers de la vente à distance. C'est
cette expertise qui l'a conduite, plus récemment, vers la formation professionnelle :
d'abord pour ses propres téléprospecteurs, puis, depuis environ un an, sous la forme d'un
organisme de formation externe préparant des apprenants à des titres professionnels reconnus
par l'État, comme le titre de Conseiller Relation Client à Distance ou celui d'Employé
Commercial. C'est cette activité de formation qui constitue le cœur de ce mémoire.

Ce changement d'échelle a fait apparaître une contrainte forte. Pour une petite structure,
faire appel à des formateurs en direct pour assurer les cours en ligne pèse lourdement sur
des marges déjà limitées. Une première réponse a consisté à diffuser des cours audio préparés
à l'avance plutôt qu'animés en direct ; mais ce modèle restait dépendant d'un travail humain
important, et un cours préenregistré ne permet pas de répondre aux questions des élèves.
L'idée s'est donc imposée de concevoir une plateforme de formation réellement autonome,
capable à la fois de produire et de diffuser les contenus, et d'accompagner les apprenants.

\subsection*{Problématique, objectifs et contributions}

Le projet étudié dans ce mémoire est précisément cette plateforme de formation autonome
alimentée par l'intelligence artificielle. Elle ne doit pas seulement diffuser des fichiers
audio, mais générer les contenus, les transformer en audio, produire des supports et des
exercices, et accompagner les élèves grâce à un assistant s'appuyant sur les documents réels
de la formation. La problématique peut alors être formulée ainsi :

\begin{quote}
\emph{Comment concevoir une plateforme de formation autonome alimentée par
l'intelligence artificielle, capable de produire, diffuser et accompagner des contenus
pédagogiques, tout en garantissant leur qualité, leur traçabilité et leur pertinence pour
les apprenants ?}
\end{quote}

Cette problématique se décline en quatre objectifs : automatiser la production des cours, en
passant d'un référentiel officiel à des journées pédagogiques complètes ; transformer ces
contenus en audios exploitables, à un coût inférieur à une production humaine répétée ;
enrichir la formation par des supports visuels et des exercices générés automatiquement ;
et permettre aux apprenants d'interroger un assistant capable de répondre à partir du contenu
réel du cours, et non comme un chatbot généraliste.

Pour y répondre, le travail apporte deux contributions principales. La première est une
\emph{pipeline} de génération de contenus pédagogiques, qui a évolué d'une génération
directe de texte vers une chaîne structurée, fondée sur un plan pédagogique explicite, une
génération section par section, des contrôles qualité ciblés et la conservation d'artefacts
intermédiaires garantissant la traçabilité. La seconde est un assistant pédagogique de type
RAG (\emph{Retrieval-Augmented Generation}), qui ancre ses réponses dans les documents
indexés de la formation pour dépasser les limites d'un modèle de langage utilisé seul.

La méthode suivie est itérative : chaque solution essayée (intervention humaine,
enregistrement différé, synthèse vocale, génération directe par modèle de langage, puis
pipeline structurée) a révélé une nouvelle contrainte — coût, qualité audio, longueur des
textes, cohérence pédagogique, conformité, accompagnement, évaluation. L'enjeu n'est pas
seulement d'ajouter de l'intelligence artificielle à une plateforme existante, mais de
l'encadrer pour la rendre fiable, contrôlable et mesurable. Les résultats sont donc discutés à
l'aide d'indicateurs concrets : réduction des interventions humaines, coût et temps de
production des contenus, et, pour l'assistant, pertinence des passages retrouvés et fidélité des
réponses produites.

\subsection*{Structure du mémoire}

Ce mémoire est structuré en trois parties. La première, \emph{Position du problème}, présente
le contexte de l'entreprise, les limites du fonctionnement initial et la problématique de la
formation autonome. La deuxième, \emph{État de l'art}, étudie les approches existantes :
plateformes de formation en ligne, synthèse vocale, génération de contenus par IA et systèmes
de questions-réponses augmentés par récupération d'information. La troisième, \emph{Étude
originale}, présente la solution développée, ses choix d'architecture, les difficultés
rencontrées, les arbitrages réalisés et les méthodes retenues pour en évaluer les apports.
